\documentclass[12pt, a4paper]{article}

\usepackage[T1]{fontenc}
\usepackage[utf8]{inputenc}
\usepackage[ngerman]{babel}
\usepackage{geometry}
\usepackage{hyperref}
\usepackage{listings}
\usepackage{xcolor}
\usepackage{booktabs}
\usepackage{parskip}
\usepackage{microtype}
\usepackage{fancyhdr}
\usepackage{titlesec}
\usepackage{enumitem}

\usepackage{helvet}
\renewcommand{\familydefault}{\sfdefault}

\geometry{left=2.5cm, right=2.5cm, top=2.5cm, bottom=2.5cm}

\definecolor{dhbwred}{RGB}{226, 0, 26}
\definecolor{codebg}{RGB}{245, 245, 245}
\definecolor{codeframe}{RGB}{200, 200, 200}
\definecolor{keyword}{RGB}{0, 0, 180}
\definecolor{comment}{RGB}{100, 100, 100}

\lstdefinestyle{base}{
  backgroundcolor=\color{codebg}, frame=single, rulecolor=\color{codeframe},
  basicstyle=\ttfamily\footnotesize, keywordstyle=\color{keyword}\bfseries,
  commentstyle=\color{comment}\itshape,
  showstringspaces=false, breaklines=true, captionpos=b,
  numbers=left, numberstyle=\tiny\color{comment}, numbersep=6pt, tabsize=4,
}
\lstdefinestyle{sql}{style=base, language=SQL}
\lstdefinestyle{bash}{style=base, language=bash}
\lstset{style=base}

\pagestyle{fancy}
\fancyhf{}
\fancyhead[L]{\small Full Stack Handwerk -- Projektbericht}
\fancyhead[R]{\small DHBW Ravensburg}
\fancyfoot[C]{\thepage}
\renewcommand{\headrulewidth}{0.4pt}

\hypersetup{colorlinks=true, linkcolor=black, urlcolor=dhbwred}

\titleformat{\section}{\large\bfseries}{\thesection}{1em}{}[\color{dhbwred}\titlerule]
\titleformat{\subsection}{\normalsize\bfseries}{\thesubsection}{1em}{}

% Hilfsmakro fuer Repo-Links
\newcommand{\repolink}[2]{\href{https://github.com/alinawagner1805-crypto/Full-Stack-Handwerk/#1}{#2}}

\begin{document}

% ── Titelseite ────────────────────────────────────────────────────────────────
\begin{titlepage}
  \centering
  \vspace*{1.5cm}
  {\color{dhbwred}\rule{\linewidth}{1.5pt}}\vspace{0.5cm}

  {\Huge\bfseries Projektbericht\\[0.4em]\large Full Stack Handwerk}

  \vspace{0.3cm}
  {\color{dhbwred}\rule{\linewidth}{1.5pt}}\vspace{1.5cm}

  {\LARGE\bfseries GET /stats -- Prediction Analytics Endpoint}\\[0.8em]
  {\large Erweiterung von Block~6: Datenbanken \& SQL}

  \vspace{2.5cm}
  \begin{tabular}{ll}
    \textbf{Verfasserin:}    & Alina Wagner \\[0.4em]
    \textbf{Matrikelnummer:} & -- \\[0.4em]
    \textbf{Kurs:}           & DSKI, DHBW Ravensburg \\[0.4em]
    \textbf{Dozent:}         & Prof.\ Dr.-Ing.\ Mark Schutera \\[0.4em]
    \textbf{Datum:}          & \today \\
  \end{tabular}

  \vfill
  {\small Dualer Partner: ZDF \\ DHBW Ravensburg -- Studiengang Data Science \& KI}
  \vspace{1cm}
  {\color{dhbwred}\rule{\linewidth}{0.8pt}}
\end{titlepage}

% ── Inhaltsverzeichnis ────────────────────────────────────────────────────────
\tableofcontents
\thispagestyle{fancy}
\newpage

% ══════════════════════════════════════════════════════════════════════════════
\section{Einleitung und Motivation}

Im Rahmen des Kurses \emph{Full Stack Handwerk} entstand PixelWise -- eine
Webanwendung zur Erkennung handgeschriebener Ziffern mit Machine Learning.
Mit Block~6 kam eine PostgreSQL-Datenbank hinzu, die jede Klassifikation in
der Tabelle \texttt{predictions} persistiert. Der einzige lesende Endpoint
im Auslieferungszustand, \texttt{GET /results}, gibt dabei die letzten
20~Eintraege chronologisch aus. Analytische Fragen wie \textit{,,Bei welcher
Ziffer ist das Modell am unsichersten?''} oder \textit{,,Hat sich die
Konfidenz nach einem Modellwechsel verbessert?''} bleiben damit unbeantwortet.

Diese Erweiterung behebt genau das: mit einem neuen Endpoint
\texttt{GET /stats}, der die Datenbank mit fortgeschrittenen SQL-Mitteln
auswertet, und einem automatisierten Run per systemd-Timer, der den Endpoint
taeglich aufruft und Ergebnisse als JSON-Snapshots speichert -- integriert
in die bestehende Serverinfrastruktur aus Block~5.

Alle Implementierungsdetails, vollstaendige SQL-Abfragen und der gesamte
Quellcode sind im zugehoerigen GitHub-Repository verfuegbar:\\
\url{https://github.com/alinawagner1805-crypto/Full-Stack-Handwerk}

% ══════════════════════════════════════════════════════════════════════════════
\section{SQL-Konzepte und Implementierung}

Der Kurs behandelt \texttt{INSERT} und einfaches \texttt{SELECT} ueber das
SQLAlchemy-ORM. Diese Erweiterung geht in der Datenbankschicht bewusst
weiter -- eine ausfuehrliche Gegenuebertstellung aller SQL-Konzepte findet
sich unter
\repolink{blob/main/docs/sql\_konzepte.md}{\texttt{docs/sql\_konzepte.md}}.

\subsection{Aggregationen: GROUP BY, COUNT, AVG und HAVING}

Die Kernabfrage gruppiert alle Predictions nach Ziffer und berechnet
Haeufigkeit sowie Durchschnittskonfidenz. Mit \texttt{HAVING} lassen sich
anschliessend Gruppen nach ihrem Aggregatwert filtern -- semantisch
verschieden von \texttt{WHERE}, das einzelne Zeilen \emph{vor} der
Aggregation entfernt und damit den Durchschnitt verfaelscht:

\begin{lstlisting}[style=sql, caption={Aggregation mit HAVING-Filter}]
SELECT prediction,
       COUNT(*)                           AS anzahl,
       ROUND(AVG(confidence)::numeric, 4) AS avg_konfidenz
FROM predictions
GROUP BY prediction
HAVING AVG(confidence) < 0.80
ORDER BY avg_konfidenz ASC;
\end{lstlisting}

Im Endpoint als optionaler Parameter: \texttt{GET /stats?confidence\_below=0.80}.

\subsection{Zeitreihe, Window Function und Filterparameter}

Fuer die tagesweise Aufschluesselung wird \texttt{created\_at} per
\texttt{CAST AS DATE} auf den reinen Datumsteil reduziert und als
\texttt{GROUP BY}-Schluessel verwendet. Ergaenzend berechnet eine
Window Function (\texttt{RANK() OVER (PARTITION BY prediction ORDER BY
confidence DESC)}) den Rang jeder Prediction innerhalb ihrer Zifferngruppe,
ohne die Zeilen zu kollabieren -- das ist der wesentliche Unterschied zu
\texttt{GROUP BY}, bei dem Originalzeilen verloren gehen.

Alle Filterparameter sind optional und kombinierbar:

\begin{center}
\begin{tabular}{lll}
  \toprule
  \textbf{Parameter} & \textbf{Beispiel} & \textbf{Wirkung} \\
  \midrule
  \texttt{model\_version}   & \texttt{?model\_version=v1}    & Nur eine Modellversion \\
  \texttt{since}            & \texttt{?since=2026-01-01}     & Ab einem Datum \\
  \texttt{confidence\_below}& \texttt{?confidence\_below=0.8}& Schwellenwertfilter \\
  \bottomrule
\end{tabular}
\end{center}

% ══════════════════════════════════════════════════════════════════════════════
\section{Performance-Analyse mit EXPLAIN ANALYZE}

Ohne Index liest PostgreSQL die gesamte Tabelle bei jeder Abfrage zeilenweise
(\emph{Sequential Scan}). Mit einem B-Tree-Index auf \texttt{created\_at}
kann die Datenbank bei gefilterten Anfragen direkt auf relevante Bereiche
zugreifen (\emph{Index Scan}). \texttt{EXPLAIN ANALYZE} fuehrt die Abfrage
tatsaechlich aus und liefert den realen Ausfuehrungsplan:

\begin{lstlisting}[style=bash, caption={EXPLAIN ANALYZE -- vor und nach dem Index}]
-- Ohne Index (Sequential Scan):
Seq Scan on predictions  (cost=0.00..5.00 rows=200)
  Planning Time: 1.3 ms  |  Execution Time: 0.9 ms

-- Nach: CREATE INDEX ix_predictions_created_at ON predictions(created_at);
Index Scan using ix_predictions_created_at on predictions
  Planning Time: 0.4 ms  |  Execution Time: 0.2 ms
\end{lstlisting}

Die Ausfuehrungszeit sinkt von 0{,}9\,ms auf 0{,}2\,ms (Faktor~4{,}5).
Entscheidend ist, dass der Vorteil mit dem Datenvolumen skaliert: Sequential
Scan waechst linear, Index Scan logarithmisch. Den vollstaendigen Query Plan
mit Kostenschaetzungen findet sich unter
\repolink{blob/main/docs/performance.md}{\texttt{docs/performance.md}}.

% ══════════════════════════════════════════════════════════════════════════════
\section{Integration in die Infrastruktur}

\subsection{Automatisierter Run per systemd-Timer}

Ein Endpoint der nur auf Anfrage antwortet, erkennt keine Veraenderungen ueber
Zeit. Deshalb wird \texttt{GET /stats} taeglich per systemd-Timer aufgerufen
(\texttt{OnCalendar=*-*-* 02:00:00}) und das Ergebnis als datierte JSON-Datei
unter \texttt{/opt/pixelwise/stats\_snapshots/} gespeichert.
\texttt{Persistent=true} stellt sicher, dass verpasste Runs nach einem
Serverneustart automatisch nachgeholt werden. Der vollstaendige Aufbau mit
Service-Unit, Timer-Unit und Shell-Skript ist dokumentiert unter
\repolink{blob/main/docs/monitoring.md}{\texttt{docs/monitoring.md}}.

\subsection{Baseline-Vergleich und Messbarkeit}

Die gespeicherten Snapshots ermoeglichen einen direkten Vorher-Nachher-
Vergleich ueber Modellversionen hinweg. Das Skript \texttt{compare\_stats.py}
liest zwei JSON-Dateien und gibt Konfidenz-Deltas je Ziffer aus:

\begin{lstlisting}[style=bash, caption={Baseline-Vergleich: v1 vs. v2}]
Gesamt-Predictions:  200 -> 487
Avg Konfidenz:    0.8741 -> 0.9012  (+0.0271)

Ziffer     Vorher    Nachher      Delta
6          0.7854     0.8231    +0.0377
9          0.7621     0.8105    +0.0484
\end{lstlisting}

Der gesamte Ablauf -- vom Einspielen der Testdaten bis zum Vergleich -- laesst
sich mit \texttt{demo/live\_demo.sh} in einem Schritt reproduzieren.

% ══════════════════════════════════════════════════════════════════════════════
\section{Probleme, Ergebnis und Ausblick}

\subsection{Aufgetretene Probleme}

Drei Probleme traten waehrend der Implementierung auf. \texttt{func.avg()}
liefert bei leerer Tabelle \texttt{None} -- abgefangen durch eine explizite
Nullpruefung vor dem \texttt{float()}-Cast. \texttt{CAST AS DATE} verhaelt
sich in SQLite und PostgreSQL unterschiedlich; generisches \texttt{str(r.day)}
loest das ohne separate Codepfade. Der Konfidenzfilter wurde initial mit
\texttt{WHERE} statt \texttt{HAVING} umgesetzt -- was den Durchschnitt
verfaelscht haette, da einzelne Predictions vor der Aggregation entfernt
worden waeren. Der Fehler fiel beim direkten Vergleich mit einer psql-Abfrage
auf.

\subsection{Ergebnis}

\texttt{GET /stats} liefert Gesamtmetriken, eine Ziffernaufschluesselung
mit \texttt{HAVING}-Filter, eine Tages-Zeitreihe und eine Window-Function-
basierte Rangberechnung. Der Index auf \texttt{created\_at} ist empirisch
mit \texttt{EXPLAIN ANALYZE} belegt. Der systemd-Timer integriert den
Endpoint in die bestehende Infrastruktur und schafft durch taeglich
gespeicherte Snapshots die Grundlage fuer messbare Modellvergleiche.

\subsection{Ausblick}

Als naechster Schritt waere ein Feedback-Endpoint sinnvoll, ueber den Nutzer
falsche Predictions korrigieren koennen. Kombiniert mit der
\texttt{by\_digit}-Aufschluesselung entstaende eine Konfusionsmatrix, und
die Snapshots bildeten automatisch die Datengrundlage fuer spaeteres
Modell-Retraining -- ohne zusaetzliche Infrastruktur.

% ══════════════════════════════════════════════════════════════════════════════
\section*{Quellen}
\addcontentsline{toc}{section}{Quellen}

\begin{itemize}[noitemsep]
  \item Schutera, M.\ (2026). \textit{Full Stack Handwerk -- Unfinished Lecture Notes}.
        DHBW Ravensburg. \url{https://github.com/Quillstacks/LectureMaterial}
  \item PostgreSQL Global Development Group (2024). \textit{PostgreSQL 16 -- EXPLAIN}.
        \url{https://www.postgresql.org/docs/16/sql-explain.html}
  \item PostgreSQL Global Development Group (2024). \textit{PostgreSQL 16 -- Window Functions}.
        \url{https://www.postgresql.org/docs/16/tutorial-window.html}
  \item SQLAlchemy Documentation (2024). \textit{Column Elements -- func}.
        \url{https://docs.sqlalchemy.org/en/20/core/functions.html}
  \item systemd Documentation (2024). \textit{systemd.timer}.
        \url{https://www.freedesktop.org/software/systemd/man/systemd.timer.html}
\end{itemize}

\end{document}
